\documentclass[11pt]{article}

\usepackage{amsmath,amssymb}
\usepackage{xurl}
\usepackage[colorlinks=true,linkcolor=blue,citecolor=blue,urlcolor=blue]{hyperref}
\setlength{\emergencystretch}{2em}

\input{command}

\title{The Mutual-Information Bound for Matrix Product Density Operators}
\author{}
\date{2026-07-31}

\begin{document}
\maketitle
\gapnote{false-source}{resolved}

\section{The assertion}

Proposition~4.5 of
Cirac--P\'erez-Garc\'ia--Schuch--Verstraete
\cite{Cirac2016MPDO_arXiv} states that any translationally invariant matrix
product density operator of bond dimension~$D$ satisfies
\begin{align}
    I_L
        &\leq I_{L+1},
    \label{eq:cpgsv17_source_monotonicity}\\
    \lim_{L\to\infty}\lim_{N\to\infty} I_L
        &= I_\infty<\infty.
    \label{eq:cpgsv17_source_iterated_limit}
\end{align}
The Appendix derives (\ref{eq:cpgsv17_source_monotonicity}) from strong
subadditivity and obtains boundedness from
\begin{align}
    I_L
        &\leq 4\log D,
    \label{eq:cpgsv17_source_bond_bound}
\end{align}
which it attributes to Wolf--Verstraete--Hastings--Cirac
\cite{Wolf2008AreaLaws}.\footnote{Source locations:
\path{Papers/1606.00608/MPDO-22-12-17-2.tex}:795--809 and 1316--1321.}

Here an MPDO is a matrix product operator whose periodic finite-chain
operators are positive semidefinite. This definition does not require a local
purification.

\section{The cited argument}

Wolf--Verstraete--Hastings--Cirac construct a mixed projected entangled-pair
state by applying a completely positive map at each site to virtual entangled
pairs \cite{Wolf2008AreaLaws}. They purify every local map, obtaining a pure
projected entangled-pair state with an additional local physical system.
Monotonicity of mutual information under the local ancillary traces then gives
\begin{align}
    I(A:B)_{\rho}
        &\leq I(A:B)_{\ketbra{\Psi}{\Psi}},
    \label{eq:cpgsv17_local_trace_data_processing}\\
    I(A:B)_{\ketbra{\Psi}{\Psi}}
        &= 2S(A)_{\ketbra{\Psi}{\Psi}},
    \label{eq:cpgsv17_pure_state_mutual_information}\\
    2S(A)_{\ketbra{\Psi}{\Psi}}
        &\leq 2\lvert\partial A\rvert\log D.
    \label{eq:cpgsv17_pure_state_boundary_bound}
\end{align}
For an interval in a periodic one-dimensional chain,
$\lvert\partial A\rvert=2$, so
(\ref{eq:cpgsv17_local_trace_data_processing})--%
(\ref{eq:cpgsv17_pure_state_boundary_bound}) imply
\begin{align}
    I(A:B)_{\rho}
        &\leq 4\log D.
    \label{eq:cpgsv17_locally_purified_bound}
\end{align}

In this argument, $D$ is the dimension of each virtual edge system in the
locally completely positive construction. It remains the virtual bond
dimension after purification of the local maps. It is not the operator bond
dimension of an arbitrary MPO representation. If the density operator of a
locally purified state is written as an MPO by pairing the ket and bra virtual
indices, its natural operator bond dimension is at most $D^2$. The two uses of
$D$ therefore refer to different representations and cannot be identified
without an additional positive factorization.

The cited proof applies to a locally purified tensor. It does not show that an
arbitrary family of positive periodic matrix product operators admits such a
purification. Cirac--P\'erez-Garc\'ia--Schuch--Verstraete explicitly warn
immediately before Definition~4.3 that an MPDO need not possess the displayed
local purification; after Theorem~4.4, they repeat that the purification
definition does not apply to every MPDO \cite{Cirac2016MPDO_arXiv}.
\footnote{Source locations:
\path{Papers/1606.00608/MPDO-22-12-17-2.tex}:744 and 786.}

Positivity of all periodic finite-chain operators does not supply the missing
factorization. De las Cuevas--Schuch--P\'erez-Garc\'ia--Cirac prove that, for
finite multipartite states, the bond dimension $D'$ of a matrix product
purification cannot be bounded in terms of the MPO bond dimension $D$ alone
\cite{DeLasCuevas2013Purifications}. More strongly, de las
Cuevas--Cubitt--Cirac--Wolf--P\'erez-Garc\'ia construct translationally
invariant MPDOs that are positive at every chain length but admit no
translationally invariant purification valid at every chain length; their
construction already applies to classical states
\cite{DeLasCuevas2016Fundamental}. These results prevent the cited proof from
being extended to every MPDO by choosing a bond-controlled local purification.
They do not, by themselves, disprove the mutual-information estimate.

\section{Why the marginal-rank argument does not repair the proof}

Cutting a periodic MPO at the two boundary bonds gives an operator-Schmidt
decomposition with at most $D^2$ summands. This decomposition does not imply
that either reduced density matrix has rank at most $D^2$.

Let $q$ be a full-rank one-site density matrix and consider the bond-dimension
one tensor with scalar entries
\begin{align}
    M^{ij}
        &= q_{ij}.
    \label{eq:cpgsv17_product_tensor_entries}
\end{align}
The $N$-site operator and its reduction to $L$ sites are
\begin{align}
    \rho^{(N)}
        &= q^{\otimes N},
    \label{eq:cpgsv17_product_chain_operator}\\
    \rho_L
        &= q^{\otimes L}.
    \label{eq:cpgsv17_product_reduced_operator}
\end{align}
Consequently,
\begin{align}
    \operatorname{rank}(q^{\otimes L})
        &= \operatorname{rank}(q)^L
        > 1
        = D^2
    \label{eq:cpgsv17_product_marginal_rank}
\end{align}
for $L>0$, although the mutual information is zero. The cancellation of the
extensive marginal entropies in
\begin{align}
    I_L
        &= S_L+S_{N-L}-S_N
    \label{eq:cpgsv17_mutual_information_entropy_difference}
\end{align}
is therefore essential. Bounding the two marginal entropies separately by
$2\log D$ is invalid for mixed MPOs.

The operator-Schmidt decomposition alone supplies no marginal-rank bound. A
different argument, based on the supported-marginal quantum channel and a
weighted Hilbert--Schmidt contraction, controls the mutual information
directly. The next section gives this argument; it is not the proof cited by
Cirac--P\'erez-Garc\'ia--Schuch--Verstraete
\cite{Cirac2016MPDO_arXiv}.

\section{Mutual information and operator-Schmidt rank}

Let $\rho_{AB}$ be a finite-dimensional density operator, and define its
operator-Schmidt rank by
\begin{align}
    r
        &= \operatorname{OSR}(\rho_{AB}).
    \label{eq:cpgsv17_operator_schmidt_rank_definition}
\end{align}
Compress both tensor factors to the supports of the marginals $\rho_A$ and
$\rho_B$. Positivity implies that $\rho_{AB}$ is supported on the tensor
product of these subspaces. Compression and the inverse isometric inclusion
preserve mutual information and operator-Schmidt rank. We may therefore set
\begin{align}
    \sigma
        &= \rho_A>0,
    \label{eq:cpgsv17_faithful_first_marginal}\\
    \tau
        &= \rho_B>0.
    \label{eq:cpgsv17_faithful_second_marginal}
\end{align}

Choose an eigenbasis of $\sigma$. The supported-marginal state--channel
correspondence gives a trace-preserving completely positive map $\Phi$ such
that
\begin{align}
    \rho_{AB}
        &= (\operatorname{id}\otimes\Phi)
            (\ketbra{\Omega_\sigma}{\Omega_\sigma}),
    \label{eq:cpgsv17_supported_marginal_reconstruction}\\
    \tau
        &= \Phi(\sigma),
    \label{eq:cpgsv17_supported_marginal_output}\\
    \dim_{\mathbb C}\operatorname{ran}\Phi
        &= r,
    \label{eq:cpgsv17_supported_marginal_range_dimension}
\end{align}
where
\begin{align}
    \ket{\Omega_\sigma}
        &= \sum_i\sqrt{p_i}\ket{i,i},
    \label{eq:cpgsv17_weighted_maximally_entangled_vector}\\
    \sigma
        &= \operatorname{diag}(p_i).
    \label{eq:cpgsv17_first_marginal_diagonalization}
\end{align}
The equality in
(\ref{eq:cpgsv17_supported_marginal_range_dimension}) follows because the
invertible scaling by $\sqrt{p_ip_j}$ preserves the span of the operator
blocks.

For a positive-definite matrix $\omega$, define
\begin{align}
    \Gamma_\omega^z(X)
        &= \omega^{z/2}X\omega^{z/2},
    \label{eq:cpgsv17_modular_congruence}
\end{align}
and set
\begin{align}
    L
        &= \Gamma_\tau^{-1/2}
            \circ\Phi
            \circ\Gamma_\sigma^{1/2}.
    \label{eq:cpgsv17_weighted_hilbert_schmidt_map}
\end{align}
Equation~(18) in the proof of Theorem~6 of Ref.~\cite{Beigi2013Sandwiched}
implies at exponent two that $L$ is a Hilbert--Schmidt contraction. Every
singular value of $L$ is therefore at most one. Since the two modular
congruences are invertible,
\begin{align}
    \operatorname{rank}L
        &= \operatorname{rank}\Phi
        = r.
    \label{eq:cpgsv17_weighted_map_rank}
\end{align}

The order-two sandwiched whitening of $\rho_{AB}$ is the unnormalized Choi
matrix of $L$:
\begin{align}
    W
        &=
        (\sigma^T\otimes\tau)^{-1/4}
        \rho_{AB}
        (\sigma^T\otimes\tau)^{-1/4},
    \label{eq:cpgsv17_whitened_state_definition}\\
    W
        &= \sum_{i,j}E_{ij}\otimes L(E_{ij}).
    \label{eq:cpgsv17_whitened_state_choi_form}
\end{align}
The matrix units form a Hilbert--Schmidt orthonormal basis, so
\begin{align}
    \tr(W^2)
        &= \sum_{i,j}\lVert L(E_{ij})\rVert_2^2,
    \label{eq:cpgsv17_choi_parseval}\\
    \sum_{i,j}\lVert L(E_{ij})\rVert_2^2
        &= \sum_k s_k(L)^2,
    \label{eq:cpgsv17_singular_value_sum}\\
    \sum_k s_k(L)^2
        &\leq r.
    \label{eq:cpgsv17_whitened_trace_rank_bound}
\end{align}

For positive-semidefinite $\rho$ and $\omega$ satisfying
$\ker\omega\subseteq\ker\rho$, define the order-two sandwiched functional by
\begin{align}
    Q_2(\rho,\omega)
        &=
        \operatorname{Re}\tr\!\left[\left(\omega^{-1/4}\rho\omega^{-1/4}\right)^2\right],
    \label{eq:cpgsv17_order_two_sandwiched_functional}
\end{align}
where the negative power vanishes on the kernel of $\omega$. A direct endpoint
comparison gives
\begin{align}
    D(\rho\Vert\omega)
        &\leq \log Q_2(\rho,\omega)
    \label{eq:cpgsv17_relative_entropy_order_two_comparison}
\end{align}
for trace-one $\rho,\omega\geq0$ satisfying the same support condition.

To prove the faithful-reference case, define
\begin{align}
    R
        &= \sqrt{\rho},
    \label{eq:cpgsv17_relative_modular_square_root}\\
    T
        &= \omega^{-1/2},
    \label{eq:cpgsv17_relative_modular_inverse_weight}\\
    \Delta
        &= R^T\otimes T,
    \label{eq:cpgsv17_relative_modular_operator}\\
    v
        &= \operatorname{vec}(R).
    \label{eq:cpgsv17_relative_modular_vector}
\end{align}
Under column stacking,
\begin{align}
    (B\otimes A)\operatorname{vec}(X)
        &= \operatorname{vec}(AXB^T).
    \label{eq:cpgsv17_vectorization_identity}
\end{align}
The factor order in $\Delta$ therefore gives
\begin{align}
    \Delta v
        &= \operatorname{vec}(T\rho),
    \label{eq:cpgsv17_relative_modular_action}\\
    m
        &= \operatorname{Re}\langle v,\Delta v\rangle,
    \label{eq:cpgsv17_relative_modular_half_moment}\\
    m
        &= \langle R,RTR\rangle_{\mathrm{HS}}.
    \label{eq:cpgsv17_half_moment_hilbert_schmidt_form}
\end{align}
The support projection of $\Delta$ fixes $v$. For positive-semidefinite
operators $A$ and $B$, the support-sensitive tensor logarithm is
\begin{align}
    \log(A\otimes B)
        &=
        (\log A)\otimes P_{\operatorname{supp}B}
        +
        P_{\operatorname{supp}A}\otimes(\log B).
    \label{eq:cpgsv17_tensor_logarithm_support_formula}
\end{align}
The support projections in
(\ref{eq:cpgsv17_tensor_logarithm_support_formula}) are essential when an
eigenvalue vanishes. With the convention $\log 0=0$, replacing them by
ambient identity operators makes the formula false. Using
\begin{align}
    \log\sqrt{\rho}
        &= \frac12\log\rho,
    \label{eq:cpgsv17_logarithm_square_root}\\
    \log(\omega^{-1/2})
        &= -\frac12\log\omega,
    \label{eq:cpgsv17_logarithm_inverse_square_root}
\end{align}
together with transpose covariance gives
\begin{align}
    \operatorname{Re}\langle v,(\log\Delta)v\rangle
        &= \frac12D(\rho\Vert\omega).
    \label{eq:cpgsv17_log_modular_relative_entropy}
\end{align}
Support-aware Jensen and Hilbert--Schmidt Cauchy--Schwarz yield
\begin{align}
    \frac12D(\rho\Vert\omega)
        &\leq \log m,
    \label{eq:cpgsv17_support_jensen_endpoint}\\
    m^2
        &\leq
        \lVert R\rVert_2^2
        \lVert RTR\rVert_2^2,
    \label{eq:cpgsv17_half_moment_cauchy_schwarz}\\
    \lVert R\rVert_2^2
        \lVert RTR\rVert_2^2
        &= Q_2(\rho,\omega).
    \label{eq:cpgsv17_order_two_functional_hilbert_schmidt}
\end{align}
Equations~(\ref{eq:cpgsv17_support_jensen_endpoint})--%
(\ref{eq:cpgsv17_order_two_functional_hilbert_schmidt}) prove
(\ref{eq:cpgsv17_relative_entropy_order_two_comparison}) for faithful
$\omega$.

The auxiliary divergence and vector-state Jensen mechanism are adapted from
Definition~5 and Lemma~19, as used in the proof of Theorem~7, of
Ref.~\cite{MullerLennert2013Renyi}. Only the direct order-two endpoint and the
endpoint Cauchy--Schwarz estimate are used here; this shortened proof is not
Beigi's interpolation argument \cite{Beigi2013Sandwiched}. Compression to the
support of $\omega$ proves the singular-reference statement. Applying
(\ref{eq:cpgsv17_relative_entropy_order_two_comparison}) with
$\omega=\rho_A\otimes\rho_B$ gives
\begin{align}
    I(A:B)_\rho
        &= D(\rho_{AB}\Vert\rho_A\otimes\rho_B),
    \label{eq:cpgsv17_mutual_information_relative_entropy}\\
    D(\rho_{AB}\Vert\rho_A\otimes\rho_B)
        &\leq \log\tr(W^2),
    \label{eq:cpgsv17_mutual_information_whitening_bound}\\
    \log\tr(W^2)
        &\leq \log r.
    \label{eq:cpgsv17_mutual_information_osr_bound}
\end{align}
Thus every finite-dimensional bipartite density operator satisfies the sharp
coefficient-one estimate
\begin{align}
    I(A:B)_\rho
        &\leq \log\operatorname{OSR}(\rho_{AB}).
    \label{eq:cpgsv17_coefficient_one_osr_bound}
\end{align}
The uniform perfectly correlated state on $r$ classical symbols attains
equality because its operator-Schmidt rank is $r$ and its mutual information
is $\log r$.

The supported-marginal state--channel correspondence and its exact range
dimension are formalized.\footnote{Formal declarations:
\leanid{Matrix.supportedMarginalMap_isKrausCPTP},
\leanid{Matrix.supportedMarginalReconstruction_eq}, and
\leanid{Matrix.finrank_range_supportedMarginalMap} in
\doclink{TNLean/Channel/SupportedMarginalChannel}.}
The Choi-matrix estimate is also formalized.\footnote{Formal declarations:
\leanid{Matrix.rectangularChoi_frobenius_parseval} and
\leanid{Matrix.rectangularChoi_frobenius_sq_le_finrank_range} in
\doclink{TNLean/Algebra/RectangularChoi}.}
The exponent-two weighted Hilbert--Schmidt contraction is formalized for both
positive-definite and arbitrary positive-semidefinite output weights, with the
negative quarter power restricted to the output support.\footnote{Formal
declarations:
\leanid{Matrix.weightedHilbertSchmidtMap_norm_le},
\leanid{Matrix.weightedHilbertSchmidtMap_isHilbertSchmidtContraction},
\leanid{Matrix.supportedWeightedHilbertSchmidtMap_norm_le}, and
\leanid{Matrix.supportedWeightedHilbertSchmidtMap_isHilbertSchmidtContraction}
in \doclink{TNLean/Channel/WeightedHilbertSchmidt}.}
The full-support eigenbasis calculation is formalized: the whitening is
identified with the rectangular Choi matrix, including the input-transpose
convention; the modular congruences preserve the range dimension; and the
resulting operator is positive semidefinite, with its trace square and squared
Frobenius norm bounded by operator-Schmidt rank.\footnote{Formal declarations:
\leanid{Matrix.supportedMarginalWhitenedState_eq_rectangularChoi},
\leanid{Matrix.finrank_range_weightedHilbertSchmidtMap}, and
\leanid{Matrix.supportedMarginalWhitenedState_trace_sq_re_le_operatorSchmidtRank}
in \doclink{TNLean/Channel/WhitenedChoi}.}

The single-reference support reduction is also formalized. Let
$\rho,\omega\geq0$ satisfy $\ker\omega\subseteq\ker\rho$, and let
$V:\mathbb C^k\to\mathbb C^D$ satisfy
\begin{align}
    V^\dagger V
        &= I_k,
    \label{eq:cpgsv17_support_isometry_left_inverse}\\
    VV^\dagger
        &= P_{\operatorname{supp}\omega}.
    \label{eq:cpgsv17_support_isometry_range_projection}
\end{align}
Then
\begin{align}
    D(\rho\Vert\omega)
        &=
        D(V^\dagger\rho V\Vert V^\dagger\omega V),
    \label{eq:cpgsv17_relative_entropy_support_compression}\\
    Q_2(\rho,\omega)
        &=
        Q_2(V^\dagger\rho V,V^\dagger\omega V).
    \label{eq:cpgsv17_order_two_support_compression}
\end{align}
The logarithm and negative quarter power vanish on the zero eigenspace. The
compressed reference is positive definite. If $\rho$ and $\omega$ have trace
one, isometric trace preservation retains the same normalization after
compression. The singular-reference inequality therefore reduces to
\begin{align}
    D(A\Vert B)
        &\leq \log Q_2(A,B),
    \label{eq:cpgsv17_faithful_relative_entropy_comparison}\\
    A
        &\geq0,
    \qquad
    B>0,
    \qquad
    \tr A=1.
    \label{eq:cpgsv17_faithful_comparison_hypotheses}
\end{align}
No normalization of $B$ is required.

The rectangular isometric-conjugation identities are formalized.
\footnote{Formal declarations:
\leanid{Matrix.isometryConjNonUnitalStarAlgHom},
\leanid{Matrix.isometryConjNonUnitalStarAlgHom_apply},
\leanid{Matrix.cfc_conj_isometry_of_zero},
\leanid{Matrix.trace_mul_mul_conjTranspose_of_conjTranspose_mul_eq_one}, and
\leanid{Matrix.conjTranspose_mul_mul_mul_conjTranspose_mul_of_isometry}
in \doclink{TNLean/Analysis/IsometricCompression}.}
Faithful compression of a positive-semidefinite matrix onto its support is
formalized by
\leanid{Matrix.PosSemidef.compression_on_support_posDef} in
\doclink{TNLean/Analysis/SupportCompression}.
The order-two functional is
\leanid{TNLean.sandwichedRenyiTwoTrace} in
\doclink{TNLean/Analysis/SandwichedRenyiTwo}.
The reconstruction and entropy identities are formalized separately.
\footnote{Formal declarations:
\leanid{Matrix.PosSemidef.eq_expansion_compression_of_kernel_le_of_mul_conjTranspose_eq_supportProj},
\leanid{TNLean.quantumRelativeEntropy_support_compression},
\leanid{TNLean.sandwichedRenyiTwoTrace_support_compression}, and
\leanid{TNLean.quantumRelativeEntropy_le_log_sandwichedRenyiTwoTrace_of_faithful}
in \doclink{TNLean/Analysis/SupportCompressedEntropy}.}

A support-aware vector-state Jensen inequality for the logarithm is also
formalized. If $A\geq0$, $\langle v,v\rangle=1$, and the support projection of
$A$ fixes $v$, then
\begin{align}
    \operatorname{Re}\langle v,(\log A)v\rangle
        &\leq
        \log\!\left(\operatorname{Re}\langle v,Av\rangle\right).
    \label{eq:cpgsv17_support_log_jensen}
\end{align}
In an eigenbasis of $A$, the support condition makes every weight at a zero
eigenvalue vanish. Scalar logarithmic concavity is therefore applied only to
positive eigenvalues; neither concavity at zero nor positive definiteness of
$A$ is assumed. The proof uses shared spectral-weight formulas for quadratic
forms of a Hermitian matrix and its functional calculus.
\footnote{Formal declarations:
\leanid{Matrix.PosSemidef.re_dotProduct_cfc_log_mulVec_le_log} in
\doclink{TNLean/Analysis/SupportLogJensen}, and
\leanid{Matrix.IsHermitian.spectralWeight},
\leanid{Matrix.IsHermitian.spectralWeight_nonneg},
\leanid{Matrix.IsHermitian.sum_spectralWeight},
\leanid{Matrix.IsHermitian.re_dotProduct_cfc_mulVec_eq_sum}, and
\leanid{Matrix.IsHermitian.re_dotProduct_mulVec_eq_sum} in
\doclink{TNLean/Analysis/SpectralQuadraticForm}.}
This auxiliary analytic result is not a theorem of
Cirac--P\'erez-Garc\'ia--Schuch--Verstraete
\cite{Cirac2016MPDO_arXiv}. It does not alter the source-coverage count or
complete Proposition~4.5.

The reusable functional-calculus identities and the entropy comparison are
formalized. Functional calculus commutes with transpose on Hermitian matrices.
In particular, the logarithm and support projection commute with transpose,
and every real power and the positive square root commute with transpose on
positive-semidefinite matrices. These identities include singular matrices and
zero-dimensional matrix algebras. The tensor-product logarithm retains both
support projections; the logarithm of a positive-semidefinite square root is
half the logarithm; and the logarithm of a real power of a positive-definite
matrix is the exponent times the logarithm.

The faithful theorem assumes only $\rho\geq0$, $\omega>0$, and $\tr\rho=1$.
The singular theorem assumes $\rho,\omega\geq0$, both traces equal to one, and
$\ker\omega\subseteq\ker\rho$. These functional-calculus identities are
project-derived auxiliary results, not theorems of
Cirac--P\'erez-Garc\'ia--Schuch--Verstraete
\cite{Cirac2016MPDO_arXiv}; in particular, the tensor-product result permits
singular factors and zero-dimensional matrix algebras.
\footnote{Formal declarations:
\leanid{Matrix.cfc_transpose},
\leanid{Matrix.IsHermitian.log_transpose},
\leanid{Matrix.PosSemidef.rpow_transpose},
\leanid{Matrix.PosSemidef.sqrt_transpose},
\leanid{Matrix.IsHermitian.supportProj_transpose},
\leanid{Matrix.log_kronecker_posSemidef},
\leanid{Matrix.PosSemidef.cfc_log_sqrt},
\leanid{Matrix.PosDef.cfc_log_rpow},
\leanid{TNLean.quantumRelativeEntropy_le_log_sandwichedRenyiTwoTrace_posDef},
and
\leanid{TNLean.quantumRelativeEntropy_le_log_sandwichedRenyiTwoTrace} in
\doclink{TNLean/Analysis/CfcConjugation},
\doclink{TNLean/Analysis/CfcKronecker},
\doclink{TNLean/Analysis/MatrixSqrt},
\doclink{TNLean/Analysis/SandwichedRenyiTwo}, and
\doclink{TNLean/Analysis/SupportCompressedEntropy}.}

The one-sided output-support compression is also formalized. Faithfulness of
the input weight forces every channel output into the support of its image,
the compressed channel is trace preserving, and the support-restricted
negative quarter power agrees with the ordinary negative quarter power after
compression. The arbitrary-support estimate below uses simultaneous
compression to both marginal supports.

The unnormalized sandwiched R\'enyi trace term is
\begin{align}
    \widetilde Q_\alpha(\rho\Vert\omega)
        &=
        \operatorname{Re}\tr\!\left[\left(\omega^{(1-\alpha)/(2\alpha)}
                \rho
                \omega^{(1-\alpha)/(2\alpha)}\right)^\alpha\right].
    \label{eq:cpgsv17_sandwiched_renyi_trace}
\end{align}
For trace-one $\rho$ and $\alpha>1$, this is the quantity inside the logarithm
of the sandwiched R\'enyi divergence when $\rho,\omega\geq0$ and
$\ker\omega\subseteq\ker\rho$. This includes singular $\omega$: powers that
vanish on the kernel implement the generalized inverse on its support. In this
regime, finiteness of the divergence requires the support inclusion. If the
inclusion fails, the totalized continuous-functional-calculus value remains
finite but is not the divergence trace term.

The present application uses $1<\alpha\leq2$; the order-one case is the
Umegaki endpoint. The formalized properties are nonnegativity for
$\rho\geq0$ and $\omega>0$ and the exact identities at $\alpha=1$ and
$\alpha=2$.\footnote{Formal declarations:
\leanid{TNLean.sandwichedRenyiTrace},
\leanid{TNLean.sandwichedRenyiTrace_nonneg},
\leanid{TNLean.sandwichedRenyiTrace_one}, and
\leanid{TNLean.sandwichedRenyiTrace_two} in
\doclink{TNLean/Analysis/SandwichedRenyi}.}
These supporting identities are not a theorem of
Cirac--P\'erez-Garc\'ia--Schuch--Verstraete
\cite{Cirac2016MPDO_arXiv}. They prove neither continuity nor monotonicity in
$\alpha$, interpolation between orders one and two, differentiability or a
limit at $\alpha=1$, the logarithmic sandwiched R\'enyi divergence, nor a
comparison with Umegaki relative entropy.

Simultaneous compression by the two marginal support isometries, exact
reconstruction, preservation of ordinary operator-Schmidt rank, the exact
formulas for the compressed marginals, and their positive definiteness are
formalized.\footnote{Formal declarations:
\leanid{Matrix.PosSemidef.eq_marginalSupport_expansion_compression},
\leanid{Matrix.PosSemidef.operatorSchmidtRank_marginalSupport_compression},
\leanid{Matrix.PosSemidef.partialTraceRight_marginalSupport_compression},
\leanid{Matrix.PosSemidef.partialTraceLeft_marginalSupport_compression}, and
\leanid{Matrix.PosSemidef.marginalSupport_compression_marginals_posDef} in
\doclink{TNLean/Channel/MarginalSupportAbsorption}.}
The mathematical tasks tracked in
\href{https://github.com/LionSR/TNLean/issues/5209}{\#5209} and
\href{https://github.com/LionSR/TNLean/issues/5242}{\#5242} are complete. The
faithful and singular relative-entropy comparisons are proved through the
relative-modular half-moment, support-aware Jensen inequality, and
Hilbert--Schmidt Cauchy--Schwarz, without a general weighted interpolation
theorem or an order-one limit.

The faithful marginal-whitening theorem diagonalizes the first marginal, so
the input transpose from the Choi convention agrees with that marginal, and
proves
\begin{align}
    Q_2(\rho_{AB},\rho_A\otimes\rho_B)
        &\leq \operatorname{OSR}(\rho_{AB})
    \label{eq:cpgsv17_faithful_product_marginal_whitening}
\end{align}
when both marginals are positive definite. The arbitrary-support theorem
compresses simultaneously to the two marginal supports, applies
(\ref{eq:cpgsv17_faithful_product_marginal_whitening}), and returns through
exact preservation of $Q_2$ and ordinary operator-Schmidt rank. It includes
zero-dimensional support coordinates and requires no trace normalization.
\footnote{Formal declarations:
\leanid{TNLean.sandwichedRenyiTwoTrace_product_marginals_le_operatorSchmidtRank_of_marginals_posDef}
in \doclink{TNLean/Channel/FaithfulMarginalWhitenedChoi}, and
\leanid{TNLean.sandwichedRenyiTwoTrace_product_marginals_le_operatorSchmidtRank}
in \doclink{TNLean/Channel/MarginalSupportWhitenedChoi}.}
Thus the mathematical work tracked by
\href{https://github.com/LionSR/TNLean/issues/5211}{\#5211} is complete.

The unrestricted coefficient-one inequality
(\ref{eq:cpgsv17_coefficient_one_osr_bound}) is also formalized for every
positive-semidefinite bipartite operator of trace one, using ordinary
operator-Schmidt rank and without positive-dimension or faithful-marginal
assumptions. The proof identifies mutual information with relative entropy
against the product of the marginals, applies the singular-reference
relative-entropy comparison and the order-two whitening bound, and treats the
zero value of the totalized order-two functional separately before using
monotonicity of the logarithm.\footnote{Formal declarations:
\leanid{Matrix.operatorSchmidtRank_pos_of_ne_zero},
\leanid{Matrix.PosSemidef.productMarginals_kernel_le},
\leanid{TNLean.sandwichedRenyiTwoTrace_submatrix_equiv}, and
\leanid{Entropy.mutualInformation_le_log_operatorSchmidtRank} in
\doclink{TNLean/Algebra/OperatorSchmidt},
\doclink{TNLean/Channel/MarginalSupportAbsorption},
\doclink{TNLean/Analysis/SandwichedRenyiTwo}, and
\doclink{TNLean/Entropy/MutualInformationOperatorSchmidt}.}
Hence the mathematical work tracked by
\href{https://github.com/LionSR/TNLean/issues/5212}{\#5212} is complete.

The finite-chain MPO application is formalized by
\leanid{MPOTensor.mutualInfoChain_le_two_log_bondDim}, and the source-facing
coefficient-four consequence is formalized by
\leanid{MPOTensor.IsMPDO.mutualInfoChain_le_four_log_bondDim}, both in
\doclink{TNLean/MPS/MPDO/MutualInfoAreaLaw}. Thus the finite-chain work tracked
by \href{https://github.com/LionSR/TNLean/issues/5213}{\#5213} is complete.
This work does not formalize the false thermodynamic-limit clause of
Proposition~4.5 of Ref.~\cite{Cirac2016MPDO_arXiv} or alter its separate
negative status.

\section{Diagonal matrix product operators}

No diagonal tensor can violate the proposed finite-chain estimate. Fix a chain
of length $N=L+R$, and define its diagonal coefficient matrix by
\begin{align}
    P_{x,y}
        &=
        \tr\!\left(M^{x_0x_0}\cdots M^{x_{L-1}x_{L-1}}
            M^{y_0y_0}\cdots M^{y_{R-1}y_{R-1}}\right),
    \label{eq:cpgsv17_diagonal_cut_matrix}
\end{align}
where $x$ and $y$ label configurations on the two consecutive blocks. Set
\begin{align}
    A_x
        &= M^{x_0x_0}\cdots M^{x_{L-1}x_{L-1}},
    \label{eq:cpgsv17_left_block_matrix}\\
    B_y
        &= M^{y_0y_0}\cdots M^{y_{R-1}y_{R-1}}.
    \label{eq:cpgsv17_right_block_matrix}
\end{align}
Then
\begin{align}
    P_{x,y}
        &= \tr(A_xB_y),
    \label{eq:cpgsv17_diagonal_cut_trace_factorization}\\
    \tr(A_xB_y)
        &= \sum_{a,b=0}^{D-1}(A_x)_{ab}(B_y)_{ba}.
    \label{eq:cpgsv17_diagonal_cut_coordinate_factorization}
\end{align}
The factorization in
(\ref{eq:cpgsv17_diagonal_cut_coordinate_factorization}) passes through a
complex vector space of dimension $D^2$, so
\begin{align}
    \operatorname{rank}_{\mathbb C}P
        &\leq D^2.
    \label{eq:cpgsv17_diagonal_cut_complex_rank}
\end{align}

Suppose the finite-chain operator is positive semidefinite and diagonal, with
positive trace
\begin{align}
    Z
        &= \tr\rho^{(N)}(M)>0.
    \label{eq:cpgsv17_diagonal_chain_mass}
\end{align}
Its joint probability matrix is
\begin{align}
    \widehat P
        &= Z^{-1}P.
    \label{eq:cpgsv17_normalized_diagonal_distribution}
\end{align}
Multiplication by the nonzero real scalar $Z^{-1}$ preserves rank. Since
$\widehat P$ is real,
\begin{align}
    \operatorname{rank}_{\mathbb R}\widehat P
        &=
        \operatorname{rank}_{\mathbb C}\widehat P,
    \label{eq:cpgsv17_real_complex_distribution_rank}\\
    \operatorname{rank}_{\mathbb C}\widehat P
        &\leq D^2.
    \label{eq:cpgsv17_normalized_diagonal_rank_bound}
\end{align}

Riazanov--Vyalyi prove for every finite joint probability matrix that
\begin{align}
    I(X:Y)
        &\leq
        \log\operatorname{rank}_{\mathbb R}\widehat P.
    \label{eq:cpgsv17_classical_information_rank_bound}
\end{align}
This is Theorem~4.1 of Ref.~\cite{Riazanov2017PSDRank}. Their rank is the
ordinary real matrix rank, not the nonnegative rank. Their proof eliminates
linearly dependent rows while preserving nonnegativity and the column
marginals. At each step it chooses an endpoint at which mutual information
does not decrease. The final distribution has at most
$\operatorname{rank}_{\mathbb R}P$ nonzero rows. Therefore every diagonal
periodic MPO satisfies
\begin{align}
    I_L^{(N)}
        &\leq 2\log D.
    \label{eq:cpgsv17_diagonal_finite_chain_bound}
\end{align}
This conclusion requires positivity and normalization only at the chain length
under consideration; positivity at other lengths is unnecessary. It gives an
independent proof of the diagonal case. The coefficient-one theorem
(\ref{eq:cpgsv17_coefficient_one_osr_bound}) rules out both classical and
genuinely quantum counterexamples to the unrestricted finite-chain estimate.

\subsection{Local correction at the row-elimination endpoints}

The displayed calculation in the proof of Theorem~4.1, Section~4.2, of
Ref.~\cite{Riazanov2017PSDRank} cancels the factor
$1+\varepsilon\alpha_i$. In the source notation,
$m_{ij}=P_{x,y}$ is a joint-probability entry,
$p_i=\sum_jm_{ij}$ is its row marginal, and
$\alpha_i=\beta_x$ is the row-perturbation coefficient. The cancelled factor
occurs in both $m_{ij}(1+\varepsilon\alpha_i)$ and
$p_i(1+\varepsilon\alpha_i)$. The proof then chooses an endpoint of the
admissible interval, where this factor vanishes for at least one row. At such a
row, the displayed quotient has the indeterminate form $0/0$, so the
cancellation does not justify the endpoint identity as printed.

This is a local proof correction; the theorem remains unchanged. Define
$h(t)=-t\log t$ for $t>0$ and $h(0)=0$, and write mutual information as
$H(X)+H(Y)-H(X,Y)$. Then
\begin{align}
    h(ab)
        &= a\,h(b)+b\,h(a).
    \label{eq:cpgsv17_zero_valid_entropy_product_identity}
\end{align}
Equation~(\ref{eq:cpgsv17_zero_valid_entropy_product_identity}) remains valid
when either factor is zero. The additional terms cancel against the
row-marginal entropy because each row marginal is the sum of its entries,
while the row relation fixes the column marginal. Mutual information is
therefore affine on the entire closed admissible interval, including both
endpoints. This removes the endpoint singularity without strengthening the
hypotheses or weakening the conclusion, so the discrepancy is a local
correction.\footnote{The corrected calculation is used in the proof of
\leanid{Entropy.classicalMutualInformation_le_log_rank} in
\doclink{TNLean/Entropy/ClassicalMutualInformation}. Its zero-valid product
identity is \leanid{Real.negMulLog_mul}.}

The rank factorization is formalized.\footnote{Formal declarations:
\leanid{MPOTensor.diagonalCutMatrix} and
\leanid{MPOTensor.diagonalCutMatrix_rank_le} in
\doclink{TNLean/MPS/MPDO/DiagonalCutRank}.}
The information-versus-rank theorem of Riazanov--Vyalyi
\cite{Riazanov2017PSDRank} is formalized by
\leanid{Entropy.classicalMutualInformation_le_log_rank} in
\doclink{TNLean/Entropy/ClassicalMutualInformation}. Real scalar normalization
preserves rank by \leanid{Matrix.rank_smul_of_mem_nonZeroDivisors}. Faithful
scalar extension preserves matrix rank by
\leanid{Matrix.rank_map_algebraMap} in
\doclink{TNLean/Algebra/MatrixRankBaseChange}. The resulting finite-chain
estimate
\leanid{MPOTensor.diagonalCut_classicalMutualInformation_le_two_log} is in
\doclink{TNLean/MPS/MPDO/DiagonalCutRank}. It assumes a real coefficient
matrix, its normalized joint distribution, and an identification of its
complex scalar extension with the diagonal cut matrix. The fixed-length
construction is formalized in
\doclink{TNLean/MPS/MPDO/DiagonalFiniteChain}.
\footnote{Formal declarations:
\leanid{MPOTensor.IsDiagonalAt},
\leanid{MPOTensor.IsDiagonal},
\leanid{MPOTensor.diagonalCutRealMatrix},
\leanid{MPOTensor.diagonalCutMass},
\leanid{MPOTensor.normalizedDiagonalCutDistribution},
\leanid{MPOTensor.diagonalCutRealMatrix_map_ofReal_of_posSemidef},
\leanid{MPOTensor.diagonalCutRealMatrix_nonneg_of_posSemidef},
\leanid{MPOTensor.diagonalCutMass_eq_trace_re},
\leanid{MPOTensor.normalizedDiagonalCutDistribution_isJointDistribution},
and
\leanid{MPOTensor.diagonalFiniteChain_classicalMutualInformation_le_two_log}.
The global diagonal-MPDO corollary is
\leanid{MPOTensor.diagonalMPDO_classicalMutualInformation_le_two_log}.}

For a positive diagonal operator at the chosen length, the real diagonal
coefficients are nonnegative, their scalar extension is the cut matrix, and
division by their positive total mass gives a joint distribution. Diagonality
identifies these normalized coefficients with the full classical finite-chain
distribution. This formal conclusion is classical and does not by itself
settle the unrestricted quantum question. The source-independent argument in
the preceding section proves the unrestricted quantum estimate; its Lean
formalization remains separate from this diagonal special case.

\section{What the rank-three separation establishes}

De las Cuevas--Schuch--P\'erez-Garc\'ia--Cirac give a rank-three separation in
Result~1 and equations~(10)--(14) of
Ref.~\cite{DeLasCuevas2013Purifications}. It is also summarized in
Section~4.3, equation~(13), and Table~2 of
Ref.~\cite{DeLasCuevas2019Separability}. Their examples are diagonal states
\begin{align}
    \rho_t
        &=
        \sum_{i,j}(M_t)_{ij}\ketbra{i,j}{i,j},
    \label{eq:cpgsv17_rank_three_diagonal_states}
\end{align}
where the nonnegative matrix $M_t$ has ordinary rank three while its positive
semidefinite rank diverges. Under their correspondence, $\rho_t$ has
operator-Schmidt rank three, while its purification rank and separable rank
diverge.

This is a statement about the size of exact factorizations, not a lower bound
on quantum mutual information. The same paper bounds the entanglement of
purification above by the logarithm of the purification rank
\cite{DeLasCuevas2013Purifications}. A diverging purification rank makes this
upper bound nonuniform but does not imply that the entanglement of purification
diverges. No lower bound on entanglement of purification or quantum mutual
information is proved for the rank-three family.

After normalizing $M_t$ to a joint probability distribution, the quantum
mutual information of the diagonal state equals the classical mutual
information of that distribution. The ordinary-rank bound in Theorem~4.1 of
Ref.~\cite{Riazanov2017PSDRank} gives
\begin{align}
    I(X:Y)
        &\leq \log\operatorname{rank}(M_t),
    \label{eq:cpgsv17_rank_three_information_rank_bound}\\
    \log\operatorname{rank}(M_t)
        &= \log 3.
    \label{eq:cpgsv17_rank_three_uniform_information_bound}
\end{align}
This theorem concerns ordinary matrix rank, not nonnegative rank. Its proof
repeatedly eliminates a linearly dependent row while preserving nonnegativity
and the column marginals, choosing an endpoint at which mutual information
does not decrease. The final joint distribution has at most
$\operatorname{rank}(M_t)$ nonzero rows, so its mutual information is at most
the entropy of a random variable with that many values.

The rank-three examples therefore have uniformly bounded mutual information.
They show that a small MPO bond dimension does not control the bond dimension
of an exact local purification; they neither prove nor disprove the
unrestricted $4\log D$ mutual-information estimate.

\section{The iterated limit is false without an aperiodicity hypothesis}

The normalized periodic states need not converge on local algebras. Consider a
tensor with physical dimension two and bond dimension three whose only nonzero
physical slices are
\begin{align}
    M^{00}
        &= \operatorname{diag}(1,0,0),
    \label{eq:cpgsv17_parity_tensor_zero_slice}\\
    M^{11}
        &= \operatorname{diag}(0,1,-1).
    \label{eq:cpgsv17_parity_tensor_one_slice}
\end{align}
For every positive chain length, contraction of the virtual trace gives
\begin{align}
    \rho^{(N)}(M)
        &=
        \ketbra{0^N}{0^N}
        +
        (1+(-1)^N)\ketbra{1^N}{1^N}.
    \label{eq:cpgsv17_parity_chain_operator}
\end{align}
The operator in (\ref{eq:cpgsv17_parity_chain_operator}) is positive
semidefinite for every $N\geq1$, including odd lengths, and
\begin{align}
    \tr\rho^{(N)}(M)
        &= 2+(-1)^N
        \neq 0.
    \label{eq:cpgsv17_parity_chain_trace}
\end{align}
This tensor is therefore an MPDO under exactly the hypotheses of
Proposition~4.5 of Ref.~\cite{Cirac2016MPDO_arXiv}.

At odd lengths, the normalized state is
\begin{align}
    \widehat\rho^{(N)}(M)
        &= \ketbra{0^N}{0^N}.
    \label{eq:cpgsv17_odd_normalized_state}
\end{align}
At positive even lengths, it is
\begin{align}
    \widehat\rho^{(N)}(M)
        &=
        \frac13\ketbra{0^N}{0^N}
        +
        \frac23\ketbra{1^N}{1^N}.
    \label{eq:cpgsv17_even_normalized_state}
\end{align}
The all-zero entry of the one-site reduced state therefore alternates between
$1$ and $1/3$, so the normalized local states have no thermodynamic limit.

The mutual information also oscillates. For fixed $L\geq1$ and every $N>L$,
the odd-length state is a product state, and hence
\begin{align}
    I_L^{(N)}
        &= 0
        \qquad\text{for odd $N$}.
    \label{eq:cpgsv17_odd_mutual_information}
\end{align}
For even $N$, both nonempty marginals and the full state have the two nonzero
eigenvalues $1/3$ and $2/3$. Define
\begin{align}
    h_2(1/3)
        &=
        -\frac13\log\frac13
        -\frac23\log\frac23
        >0.
    \label{eq:cpgsv17_binary_entropy_one_third}
\end{align}
The even-length mutual information is
\begin{align}
    I_L^{(N)}
        &=
        h_2(1/3)+h_2(1/3)-h_2(1/3),
    \label{eq:cpgsv17_even_mutual_information_entropy_difference}\\
    I_L^{(N)}
        &= h_2(1/3)
        \qquad\text{for even $N$}.
    \label{eq:cpgsv17_even_mutual_information}
\end{align}
Equations~(\ref{eq:cpgsv17_odd_mutual_information}) and
(\ref{eq:cpgsv17_even_mutual_information}) show that
$\lim_{N\to\infty}I_L^{(N)}$ does not exist. The eigenvalue $-1$ of the
physical-trace transfer operator is the peripheral sector responsible for the
parity oscillation. Primitivity or another aperiodicity hypothesis would
exclude this example, but Proposition~4.5 of
Ref.~\cite{Cirac2016MPDO_arXiv} contains no such hypothesis.

The finite-chain form, positivity, trace formula, alternating normalized
states and marginals, exact entropies and mutual informations, and
nonconvergence of both the one-site entry and mutual information are
formalized.\footnote{Formal declarations:
\leanid{MPOTensor.ThermodynamicLimitCounterexample.mpo_parityTensor_eq_diagonal},
\leanid{MPOTensor.ThermodynamicLimitCounterexample.parityTensor_isMPDO},
\leanid{MPOTensor.ThermodynamicLimitCounterexample.trace_mpo_parityTensor},
\leanid{MPOTensor.ThermodynamicLimitCounterexample.reducedBlockState_one_zero_zero_of_odd},
\leanid{MPOTensor.ThermodynamicLimitCounterexample.reducedBlockState_one_zero_zero_of_even},
\leanid{MPOTensor.ThermodynamicLimitCounterexample.reducedBlockState_parityTensor_of_even},
\leanid{MPOTensor.ThermodynamicLimitCounterexample.blockEntropy_parityTensor_of_even},
\leanid{MPOTensor.ThermodynamicLimitCounterexample.mutualInfoChain_parityTensor_of_odd},
\leanid{MPOTensor.ThermodynamicLimitCounterexample.mutualInfoChain_parityTensor_of_even},
\leanid{MPOTensor.ThermodynamicLimitCounterexample.not_exists_tendsto_reducedBlockState_one_zero_zero},
\leanid{MPOTensor.ThermodynamicLimitCounterexample.not_exists_tendsto_parityMutualInfoAfterCut},
and the source-facing conjunction
\leanid{MPOTensor.ThermodynamicLimitCounterexample.proposition45_limit_counterexample}
in \doclink{TNLean/MPS/MPDO/ThermodynamicLimitCounterexample}.}

Accordingly, the precise unconditional result retained from the cited argument
is the finite-chain monotonicity
\begin{align}
    I_L^{(N)}
        &\leq I_{L+1}^{(N)}
        \qquad
        (2L+1\leq N).
    \label{eq:cpgsv17_unconditional_finite_chain_monotonicity}
\end{align}

\section{Resolution}

The finite-chain analytic implication used by the cited proof is now
formalized.\footnote{Formal declaration:
\leanid{MPOTensor.mutualInfoChain_le_of_bipartitioned_channel_image} in
\doclink{TNLean/MPS/MPDO/LocalPurificationAreaLaw}.}
If the normalized state across a cut is the image of a normalized pure matrix
product state of bond dimension $D'$ under independent local
trace-preserving completely positive maps, then data processing and the
pure-state boundary estimate give
\begin{align}
    I_L
        &\leq 4\log D'.
    \label{eq:cpgsv17_purifying_bond_information_bound}
\end{align}

The local-purification step is also complete.\footnote{Formal declarations:
\leanid{MPOTensor.bipartitionedNormalizedMPO_eq_tensorMapBoth_blockAncillaryTraceMap},
\leanid{MPOTensor.mutualInfoChain_le_of_lpdoWitness}, and
\leanid{MPOTensor.IsLPDO.exists_mutualInfoChain_le_purifyingBond} in
\doclink{TNLean/MPS/MPDO/LocalPurificationAreaLaw}.}
For explicit matrices $A^{(i,k)}$ witnessing the locally purified form, the
two blockwise ancillary traces of the normalized purifying pure state equal the
normalized operator generated by $M$. Hence, on every nonempty chain, a
locally purified tensor with nonzero finite-chain trace satisfies
(\ref{eq:cpgsv17_purifying_bond_information_bound}), where $D'$ is the bond
dimension of the purifying matrix product state. The bond-space identification
in the locally purified form identifies $D$ with $(D')^2$ at the level of
cardinalities, but the proved estimate is stated in terms of $D'$, not the MPO
bond dimension $D$.

This result does not produce a local purification for an arbitrary positive
MPO. Cirac--P\'erez-Garc\'ia--Schuch--Verstraete warn that such a
representation need not exist \cite{Cirac2016MPDO_arXiv}, and the no-go
results of
Refs.~\cite{DeLasCuevas2013Purifications,DeLasCuevas2016Fundamental} show that
MPO positivity and bond dimension do not recover it uniformly.

The finite-chain estimate is nevertheless valid for arbitrary positive
periodic MPO families of bond dimension $D$ at every chain length for which
the periodic operator has positive trace. After normalization, opening the two
virtual boundary bonds gives
\begin{align}
    \operatorname{OSR}
        (\widehat\rho^{(N)}_{L:N-L})
        &\leq D^2.
    \label{eq:cpgsv17_mpo_cut_osr_bound}
\end{align}
Combining (\ref{eq:cpgsv17_mpo_cut_osr_bound}) with
(\ref{eq:cpgsv17_coefficient_one_osr_bound}) gives the stronger conclusion
\begin{align}
    I_L^{(N)}
        &\leq 2\log D.
    \label{eq:cpgsv17_finite_chain_two_log_bound}
\end{align}
The operator-Schmidt rank estimate is formalized.\footnote{Formal declaration:
\leanid{MPOTensor.operatorSchmidtRank_bipartitionedNormalizedMPO_le} in
\doclink{TNLean/MPS/MPDO/MutualInfoBridge}.}
The coefficient-one entropy implication is formalized by
\leanid{Entropy.mutualInformation_le_log_operatorSchmidtRank} in
\doclink{TNLean/Entropy/MutualInformationOperatorSchmidt}. The resulting
finite-chain estimate and its coefficient-four source-facing consequence are
formalized by
\leanid{MPOTensor.mutualInfoChain_le_two_log_bondDim} and
\leanid{MPOTensor.IsMPDO.mutualInfoChain_le_four_log_bondDim} in
\doclink{TNLean/MPS/MPDO/MutualInfoAreaLaw}.

The thermodynamic-limit clause of Proposition~4.5 of
Ref.~\cite{Cirac2016MPDO_arXiv} is formally refuted at its stated hypothesis
set by
\leanid{MPOTensor.ThermodynamicLimitCounterexample.proposition45_limit_counterexample}.
The finite-chain monotonicity and bond-dimension bounds remain valid as
separate results. Primitivity, aperiodicity, or another convergence condition
would exclude the parity obstruction, but such a condition is a genuine
boundary of a corrected theorem and is absent from the printed proposition.
This resolves the false-source limit clause tracked by
\href{https://github.com/LionSR/TNLean/issues/7070}{\#7070}.

\bibliographystyle{alpha}
\bibliography{references}

\end{document}
